\begin{figure}[!t]
    \centering
    \includegraphics[width=0.95\textwidth]{figs/files/pipeline.pdf}
    \caption{An illustration of the \framework.}
    \label{fig:pipeline}
    \vspace{-1em}
\end{figure}

% 把建立pip的思想說清楚。“把多agent組織成一個有效率的team，並讓他們能在真實環境中更高效的完成任務”。
% 1. 因為多agent比 單agent好已經有人研究過
% 2. 對話框架這個在 CahtDev, ChatEval, Heng Ji's 文章說過
% 3. 與環境交互 單 Agent 的文章也提過
% -- 所以應該定為在於說，提出pipeline，這pipeline能把多Agents，建立一個有效率的Team，這個team的人員可以 “依據當前的環境動態變化(替換)，去調整組織，以更好地做出決策”
% 因此我們提出AgentPip - 主要模擬一個Group在accomplish Goal時的過程，這過程需要，1.諮詢新的人.. 2.